\documentclass[11pt]{ctexart}

\usepackage[a4paper,margin=2.2cm]{geometry}
\usepackage{amsmath,amssymb}
\usepackage{booktabs}
\usepackage{xcolor}
\usepackage{tikz}
\usetikzlibrary{positioning,arrows.meta}
\usepackage{hyperref}
\hypersetup{colorlinks=true,linkcolor=blue!60!black,urlcolor=blue!60!black}
\definecolor{myred}{RGB}{192,0,0}

\title{\bfseries Project 1：用 NumPy 复现 OpenCLIP ViT-B-32 图像编码器}
\author{姓名：\underline{\hspace{4em}}\quad 学号：\underline{\hspace{6em}}}
\date{\today}

\begin{document}

\maketitle
\vspace{-2em}

\section{项目概述}

本项目要求在不依赖 PyTorch、OpenCLIP 或 GPU 的前提下，仅使用 \textbf{NumPy} 复现
OpenCLIP ViT-B-32 视觉编码器的前向计算：给定已预处理好的图像张量与预训练参数
$\theta$，输出与独立参考实现数值一致的图像特征向量。

\begin{itemize}
  \item 输入：float32 数组 $[B,3,224,224]$（B 张 RGB 图）。
  \item 输出：float32 数组 $[B,512]$，每张图一行、长度为 1 的归一化向量。
  \item 中间结构：12 层 Transformer 块，每层内部 token 序列为 $[B,50,768]$
        （1 个 CLS + 49 个 patch token）。
  \item 校验：与独立 PyTorch FP32 前向结果比对，匹配标准为
        $\lvert a-r\rvert \le \texttt{atol} + \texttt{rtol}\cdot\lvert r\rvert$。
\end{itemize}

\section{方法：数据流与本人实现}

图~\ref{fig:pipeline} 给出整个编码器的数据流动。图中 \textcolor{myred}{\textbf{红色加粗}}
标注的六个步骤为本人独立实现的部分（文件 \texttt{image\_encoder.py} 中的
\texttt{TODO-01} $\sim$ \texttt{TODO-06}）。

\begin{figure}[h]
\centering
\begin{tikzpicture}[
  node distance=6mm,
  box/.style 2 args={draw=#1, fill=#2, rounded corners=2pt,
              minimum width=46mm, minimum height=9mm, align=center, font=\small},
  arr/.style={-{Stealth[length=2.6mm]}, semithick}
]
\node[box={gray!60!black}{gray!10}] (img)  {\texttt{images [B,3,224,224]}};
\node[box={myred}{myred!6}, below=of img] (pat) {\textcolor{myred}{\textbf{TODO-01 patchify}}\\\texttt{[B,49,3072]}};
\node[box={myred}{myred!6}, below=of pat] (emb) {\textcolor{myred}{\textbf{TODO-05 embed}}\\\texttt{[B,50,768]}};
\node[box={myred}{myred!6}, below=of emb] (blk) {\textcolor{myred}{\textbf{TODO-02/03/04}}\;12$\times$块\\\texttt{[B,50,768]}};
\node[box={myred}{myred!6}, below=of blk] (out) {\textcolor{myred}{\textbf{TODO-06 CLS+投影}}\\\texttt{[B,512]}};
\draw[arr] (img) -- (pat);
\draw[arr] (pat) -- (emb);
\draw[arr] (emb) -- (blk);
\draw[arr] (blk) -- (out);
\end{tikzpicture}
\caption{图像编码器数据流。红色框为本人实现的六个 TODO。}
\label{fig:pipeline}
\end{figure}

\subsection{\textcolor{myred}{TODO-01} patchify：切成 49 个补丁}

把 $224\times224$ 图切成 $7\times7=49$ 个 $32\times32$ 补丁，每个补丁展平成
$3\times32\times32=3072$ 维。用两步 \texttt{reshape} 夹一步 \texttt{transpose} 完成：

\begin{equation*}
[B,3,224,224] \xrightarrow{\texttt{reshape}} [B,3,7,32,7,32]
\xrightarrow{\texttt{transpose}(0,2,4,1,3,5)} [B,7,7,3,32,32]
\xrightarrow{\texttt{reshape}} [B,49,3072].
\end{equation*}

\subsection{\textcolor{myred}{TODO-02} split\_heads：拆出 12 个头}

把每个 token 的 768 维特征连续切成 $12\times64$，再交换轴序：
$[B,N,768]\to[B,N,12,64]\to[B,12,N,64]$。

\subsection{\textcolor{myred}{TODO-03} 多头注意力}

核心公式 $\mathrm{Attention}(Q,K,V)=\mathrm{softmax}\!\big(QK^\top/\sqrt{d}\big)V$，
其中 $d=64$ 为每个头的维数。实现要点：$K$ 的转置只交换最后两轴
（\texttt{transpose(0,1,3,2)}），随后对 $V$ 加权求和，再合并 12 个头并做输出投影。

\subsection{\textcolor{myred}{TODO-04} 残差连接与 MLP}

每个块包含两次残差：先算注意力增量 $\mathrm{update}$ 并加回原输入
$\mathrm{residual}=x+\mathrm{update}$；再做两层 MLP（$768\to3072\to768$，中间为
GELU 激活），得到 $\mathrm{mlp\_update}$，最终 $\mathrm{return}=\mathrm{residual}+\mathrm{mlp\_update}$。

\subsection{\textcolor{myred}{TODO-05} embed：投影、CLS 与位置编码}

补丁投影 $x=\mathrm{patches}@W^\top$（无偏置），在最前拼接 CLS 行，再加位置编码。

\subsection{\textcolor{myred}{TODO-06} 输出 token}

取第 0 行 CLS，做 post-Norm，再用 $[768,512]$ 的投影矩阵投影到 512 维，最后单位化。
注意该投影矩阵已是 $[\mathrm{in},\mathrm{out}]$ 布局，\textbf{无需转置}。

\section{数值验证结果}

运行 \texttt{python check.py --stage all}，八个阶段全部 PASS，最大绝对误差如下：

\begin{table}[h]
\centering
\begin{tabular}{lcccr}
\toprule
阶段 & 形状 & 最大绝对误差 & 平均绝对误差 & 容差 \\
\midrule
preprocess & $(10,3,224,224)$ & $0.0$        & $0.0$        & atol $2\!\times\!10^{-6}$ \\
patches    & $(2,49,3072)$     & $0.0$        & $0.0$        & atol $2\!\times\!10^{-4}$ \\
tokens     & $(2,50,768)$      & $4.01\!\times\!10^{-5}$ & $1.29\!\times\!10^{-6}$ & atol $2\!\times\!10^{-4}$ \\
block0     & $(2,50,768)$      & $7.06\!\times\!10^{-6}$ & $3.25\!\times\!10^{-7}$ & atol $2\!\times\!10^{-4}$ \\
last\_tokens & $(2,50,768)$    & $4.77\!\times\!10^{-5}$ & $2.70\!\times\!10^{-6}$ & atol $2\!\times\!10^{-4}$ \\
output     & $(10,512)$        & $5.44\!\times\!10^{-7}$ & $8.31\!\times\!10^{-8}$ & atol $2\!\times\!10^{-5}$ \\
batch      & $(3,512)$         & $4.62\!\times\!10^{-7}$ & $7.63\!\times\!10^{-8}$ & atol $2\!\times\!10^{-5}$ \\
\bottomrule
\end{tabular}
\caption{各阶段与独立 PyTorch 参考实现的最大/平均绝对误差。}
\end{table}

\texttt{patches} 阶段为零误差，因为 \texttt{reshape/transpose} 只是重排标签、不动数据；
后续阶段出现 $10^{-5}\!\sim\!10^{-7}$ 量级误差，源自矩阵乘法浮点累加顺序与参考实现的
微小差异，均远小于容差。

\section{一次形状变化详解}

以 \texttt{patchify} 为例说明形状变化的思路：$[B,3,224,224]$ 中 $224$ 被拆成
$7\times32$，即把高度、宽度各分成「网格轴(7)」与「块内像素轴(32)」两级。第一次
\texttt{reshape} 只是重新贴标签、数据顺序不变；\texttt{transpose(0,2,4,1,3,5)} 把两个
网格轴提到前面、把块内轴 $(C,32,32)$ 挪到后面；最后一次 \texttt{reshape} 把
$7\times7$ 压成 49、把 $3\times32\times32$ 压成 3072。

\section{调试记录：一个修复的 bug}

初版 \texttt{split\_heads} 曾将形参 \texttt{x} 误写为大写 \texttt{X}：

\begin{verbatim}
    B,N,D=X.shape        # NameError: name 'X' is not defined
\end{verbatim}

Python 大小写敏感，\texttt{X} 与 \texttt{x} 是两个不同名字，运行时报
\texttt{NameError}。定位后统一改为小写 \texttt{x}，问题解决。此 bug 提醒：变量名大小写
自始至终必须一致。

\section{贡献与协助声明}

\begin{itemize}
  \item \textbf{个人贡献}：独立完成 \texttt{image\_encoder.py} 中全部六个 TODO 的
        实现与调试，理解并复现了 patchify、多头注意力、残差连接、MLP、CLS 输出等
        全部计算逻辑。
  \item \textbf{使用协助}：使用了 AI 助手（WorkBuddy）进行原理讲解、实现思路
        提示与数值验证辅助；所有代码最终由本人理解并亲手完成。
\end{itemize}

\section*{参考文献}

\begin{enumerate}
  \item Dosovitskiy et al., \emph{An Image is Worth 16x16 Words}, ICLR 2021.
  \item Ilharco et al., \emph{OpenCLIP}, 2021.
\end{enumerate}

\end{document}
